\chapter{Introdução}
\label{cap:introducao}

A circulação de conteúdos culturais em plataformas digitais tornou-se um fenômeno central para compreender como músicas, vídeos, estilos, linguagens e tendências alcançam visibilidade em larga escala. Em serviços de streaming musical, charts diários de popularidade e viralidade registram a trajetória de músicas que emergem, atingem picos de atenção e desaparecem gradualmente. Em plataformas de vídeo curto, por exemplo, sons, hashtags e formatos audiovisuais também se combinam em ciclos rápidos de adoção e remixagem. Em ambos os casos, a difusão cultural não ocorre apenas como uma série temporal agregada: ela é mediada por relações entre entidades distintas e por dependências temporais.

Modelos clássicos de difusão, incluindo formulações inspiradas em epidemiologia e em cascatas de influência, oferecem uma primeira aproximação para esse problema \citep{kempe2003ic,oliveira2025brasnam}. Ao representar a popularização de um conteúdo como processo de contágio, modelos como SIR e SIS permitem estimar parâmetros globais de adoção, persistência e declínio. Essa leitura é útil porque transforma curvas de popularidade em objetos matemáticos interpretáveis. Entretanto, tais modelos tendem a operar sobre populações agregadas e, por isso, têm dificuldade para representar a estrutura relacional que condiciona a difusão: artistas compartilham gêneros, músicas apresentam trajetórias semelhantes, gêneros se relacionam por coocorrência e diferentes faixas podem ocupar nichos culturais próximos.

Este projeto parte da hipótese de que redes neurais em grafos heterogêneos com componente temporal são mais adequadas para modelar esse tipo de fenômeno no contexto de streaming musical brasileiro. A heterogeneidade permite representar múltiplos tipos de nós e arestas; a dimensão temporal permite distinguir relações já observáveis em um instante das relações que só aparecem no futuro; e o aprendizado neural em grafos permite combinar atributos locais, estrutura de vizinhança e evolução temporal em uma mesma representação. Assim, a proposta desloca o foco de curvas populacionais isoladas para grafos dinâmicos nos quais a difusão cultural emerge da interação entre músicas, artistas, gêneros e trajetórias de popularidade.

A investigação concentra-se, neste momento, no domínio de música como objeto empírico principal. O projeto propõe reproduzir um baseline epidemiológico baseado em SIR sobre charts brasileiros e construir um grafo heterogêneo temporal envolvendo músicas, artistas e gêneros, a partir do dataset MGD+ \citep{seufitelli2023dsw}. Essa escolha permite comparar diretamente modelos populacionais e modelos baseados em grafos em um cenário já parcialmente consolidado, com dados de streaming musical brasileiro e métricas de viralidade e sucesso derivadas de charts diários.

O dataset combina o MGD+ \citep{seufitelli2023dsw} — que fornece metadados, atributos acústicos e redes de gênero — com séries de charts diários do Spotify Brasil (Top~200 e Viral~50) entre 2017 e 2021, restritas ao subconjunto de 1.981 músicas que aparecem em ambos os charts no período (\emph{viral}$\cap$\emph{hit}), universo herdado do trabalho replicado.

Esse dataset impõe três limitações que condicionam diretamente as afirmações que este trabalho pode fazer. Primeiro, a cobertura temporal por música é curta e concentrada: 62\% das 1.981 músicas (1.214) aparecem em um único ano civil, e a permanência em chart por par (música, chart) tem mediana de 4 semanas — o que restringe qualquer conclusão sobre dinâmica de longo prazo a uma minoria de casos. Segundo, no período reservado para teste, cerca de 95\% dos alvos de previsão correspondem a semanas em que a música está fora do chart, no valor de piso; a comparação entre modelos permanece válida porque todos são avaliados sobre as mesmas semanas, mas o RMSE agregado nesse período é dominado por acertos de ausência, e não por acertos de popularidade — por isso a leitura restrita às semanas em que a música está de fato no chart (recorte on-chart) é a leitura principal dos resultados, e não a leitura sobre o período completo. Terceiro, o \emph{split} temporal usado neste momento — treino até junho de 2020, teste a partir de dezembro de 2020 — deixa todo o período de teste dentro da pandemia de covid-19, enquanto o treino é quase inteiramente anterior a ela; qualquer diferença de desempenho entre modelos pode estar confundida com essa mudança de regime, e não apenas com a capacidade de cada modelo de capturar estrutura relacional, o que impede atribuir uma vantagem observada exclusivamente à estrutura do grafo sem uma checagem de robustez sob um período sem essa quebra de distribuição.

A pergunta central que orienta o trabalho é: em que medida a estrutura heterogênea e a dinâmica temporal explícita melhoram a modelagem de difusão de músicas em plataformas brasileiras de streaming, em comparação com baselines populacionais e GNNs estáticas? Essa pergunta se desdobra em três subquestões. Primeiro, interessa avaliar se uma rede neural em grafo heterogêneo temporal sobre relações artista--música--gênero supera o ajuste SIR em tarefas de reconstrução e predição de viralidade e sucesso. Segundo, busca-se entender quais tipos de relação contribuem mais para a previsibilidade, incluindo autoria, associação a gêneros, coocorrência entre gêneros e similaridade de trajetória entre músicas. Terceiro, investiga-se em que medida a interpretabilidade epidemiológica do baseline pode ser preservada ao se incorporar estrutura relacional, isto é, se o modelo em grafo ainda permite recuperar quantidades com significado epidemiológico, como taxas de adoção e declínio, em vez de operar apenas como caixa-preta.

Além do interesse científico dessas questões, o problema tem valor aplicado imediato. Um preditor confiável de curto horizonte para a trajetória de popularidade de uma faixa subsidiaria decisões hoje tomadas majoritariamente por inspeção manual das curvas de chart, como a priorização de faixas para curadoria de playlist, a alocação de investimento promocional e a detecção antecipada de músicas em ascensão.

As contribuições esperadas deste projeto são duas. Primeiro, construir uma representação em grafo heterogêneo temporal para o domínio musical, permitindo avaliar o ganho marginal de relações como autoria, gênero, coocorrência e similaridade de trajetória. Segundo, comparar modelos populacionais, GNNs estáticas e GNNs temporais em um mesmo pipeline experimental, produzindo uma análise sobre o papel de homofilia, estrutura relacional e temporalidade na previsibilidade da difusão musical.

Este documento de qualificação está organizado da seguinte forma. O Capítulo~\ref{cap:fundamentacao} apresenta os conceitos necessários para a proposta, incluindo redes sociais, difusão de informação, modelos epidemiológicos, grafos heterogêneos temporais e redes neurais em grafos. Os capítulos seguintes deverão detalhar os trabalhos relacionados, os dados e métodos da linha musical, os resultados parciais já obtidos e o cronograma de execução da dissertação.
